\documentclass[11pt,a4paper]{article}

\usepackage[a4paper,top=18mm,bottom=19mm,right=19mm,left=19mm]{geometry}
\usepackage{amsmath,amssymb}
\usepackage{xcolor}
\usepackage{enumitem}
\usepackage[explicit]{titlesec}
\usepackage[most]{tcolorbox}
\usepackage{fancyhdr}
\usepackage[colorlinks=true,linkcolor=blue!45!black]{hyperref}
\usepackage{xepersian}

\settextfont[
  Path={../booklet/font/},
  UprightFont={IRLotusICEE.ttf},
  BoldFont={IRLotusICEE_Bold.ttf},
  BoldItalicFont={IRLotusICEE_BoldIranic.ttf},
  ItalicFont={IRLotusICEE_Iranic.ttf},
  FontFace={b}{n}{IRLotusICEE_Bold.ttf},
  FontFace={bx}{n}{IRLotusICEE_Bold.ttf},
  FontFace={b}{it}{IRLotusICEE_BoldIranic.ttf},
  FontFace={bx}{it}{IRLotusICEE_BoldIranic.ttf},
  FontFace={m}{sl}{IRLotusICEE_Iranic.ttf},
  Scale=1.15
]{IRLotusICEE.ttf}
\setlatintextfont[
  Path={../booklet/font/},
  BoldFont={LiberationSerif-Bold.ttf},
  BoldItalicFont={LiberationSerif-BoldItalic.ttf},
  ItalicFont={LiberationSerif-Italic.ttf}
]{LiberationSerif-Regular.ttf}
\setdigitfont[Path={../booklet/font/},Scale=1.15]{IRLotusICEE.ttf}
\defpersianfont\titlefont[Path={../booklet/font/},Scale=1]{IRTitr.ttf}

\definecolor{navy}{RGB}{20,65,105}
\definecolor{greenback}{RGB}{238,248,241}
\definecolor{yellowback}{RGB}{255,248,225}

\newcommand{\eng}[1]{\lr{#1}}
\newcommand{\term}[2]{\textbf{#1} \eng{(#2)}}
\newcommand{\key}[1]{\textcolor{navy}{\textbf{#1}}}

\newtcolorbox{formula}[1]{
  enhanced,breakable,colback=greenback,colframe=green!45!black,title={#1},
  fonttitle=\bfseries,boxrule=.6pt,arc=1.5mm,
  right=3.5mm,left=3.5mm,top=2mm,bottom=2mm,
  before skip=7pt,after skip=7pt,
  boxed title style={left=2mm,right=2mm,top=.7mm,bottom=.7mm}
}
\newtcolorbox{trap}[1]{
  enhanced,breakable,colback=yellowback,colframe=orange!65!black,title={#1},
  fonttitle=\bfseries,boxrule=.6pt,arc=1.5mm,
  right=3.5mm,left=3.5mm,top=2mm,bottom=2mm,
  before skip=7pt,after skip=7pt
}

\setlist[itemize,1]{
  label=\textcolor{navy}{\large\textbullet},
  rightmargin=7mm,leftmargin=2mm,topsep=4pt,itemsep=3.5pt,parsep=1pt
}
\setlist[itemize,2]{
  label=\textcolor{navy}{--},
  rightmargin=7mm,leftmargin=2mm,topsep=2pt,itemsep=2.5pt,parsep=0pt
}
\setlist[enumerate,1]{
  label=\textcolor{navy}{\arabic*.},
  rightmargin=9mm,leftmargin=2mm,topsep=4pt,itemsep=3pt
}
\setlength{\parindent}{0pt}
\setlength{\parskip}{3pt}
\setlength{\fboxsep}{5pt}
\linespread{1.12}

\titleformat{\section}[block]
  {\normalfont\titlefont\Large\color{white}}
  {}{0pt}
  {\colorbox{navy}{\parbox{\dimexpr\textwidth-2\fboxsep\relax}{\raggedleft\strut#1\strut}}}
\titlespacing*{\section}{0pt}{13pt}{10pt}

\titleformat{\subsection}[hang]
  {\normalfont\bfseries\large\color{navy}}
  {\thesubsection}{7pt}{#1}
  [\vspace{2pt}\color{navy!35}\titlerule]
\titlespacing*{\subsection}{0pt}{11pt}{7pt}

\pagestyle{fancy}
\fancyhf{}
\fancyhead[R]{\small مرور فصل ۷: مدل‌های اسپایکی مدرن}
\fancyhead[L]{\small علوم شناختی محاسباتی}
\fancyfoot[C]{\thepage}
\setlength{\headheight}{14pt}

\begin{document}

\begin{center}
  {\titlefont\LARGE خلاصهٔ امتحانی فصل ۷}\\[3pt]
  {\Large مدل‌های اسپایکی مدرن}\\[3pt]
  {\small مرور تفکیک‌شدهٔ مفاهیم، معماری‌ها و فرمول‌ها}
\end{center}

\section{بخش اول: \eng{Spiking Transformer} و \eng{SpikeFormer}}

\subsection{از \eng{Deep SNN} به \eng{Transformer} اسپایکی}

\begin{itemize}
  \item در \eng{SNN}، داده زمان‌محور و ورودی و خروجی معمولاً اسپایک‌های صفر و یک‌اند؛ مزیت اصلی، پردازش رویدادمحور و مصرف انرژی کمتر است.
  \item در \eng{Deep SNN} نظارت‌شده، عبور گرادیان از تابع گسستهٔ اسپایک با \term{گرادیان جایگزین}{Surrogate Gradient} ممکن می‌شود.
  \item لایهٔ خطی اسپایکی، وزن را در اسپایک ضرب و سپس خروجی اسپایکی جدیدی تولید می‌کند:
  \[
  W\times\mathrm{spike}
  \]
  \item \textbf{\eng{SpikeFormer}}: ترکیب معماری \eng{Transformer} با شبکهٔ عصبی اسپایکی برای دستیابی به مدل قدرتمندتر و کم‌مصرف‌تر.
  \item سه جزء ضروری:
  \begin{enumerate}
    \item کدگذاری دادهٔ ورودی به اسپایک؛
    \item توجه اسپایکی؛
    \item بلوک \eng{MLP} اسپایکی.
  \end{enumerate}
\end{itemize}

\subsection{کدگذاری ورودی}

\begin{itemize}
  \item دادهٔ پیوسته باید به دنباله‌ای از صفر و یک در چند گام زمانی تبدیل شود.
  \item \textbf{\eng{Time Step}} با \(T\) نشان داده می‌شود؛ مثلاً \(T=4\) یعنی هر مقدار در چهار گام زمانی بازنمایی می‌شود.
  \item دو هدف کدگذاری:
  \begin{itemize}
    \item تبدیل مقدار پیوسته به اسپایک؛
    \item پخش اطلاعات در طول زمان، به‌طوری که مقدار بزرگ‌تر معمولاً نرخ اسپایک بیشتری بسازد.
  \end{itemize}
\end{itemize}

\subsection{کدگذاری تصویر و \eng{Rate Encoding}}

\begin{formula}{نرمال‌سازی پیکسل}
\[
\boxed{x_{\mathrm{norm}}=\frac{x}{255}}
\]
\end{formula}

\begin{itemize}
  \item برای \(T=4\)، تعداد تقریبی اسپایک از \(x_{\mathrm{norm}}T\) به دست می‌آید.
  \item مثال:
  \[
  [0.2,0.8,0.5,0.5]\times4=[0.8,3.2,2,2]
  \]
  \[
  \mathrm{round}(\cdot)=[1,3,2,2]
  \]
  یعنی ویژگی‌ها به‌ترتیب یک، سه، دو و دو اسپایک در کل بازه تولید می‌کنند.
  \item \textbf{\eng{Rate Encoding}}: مقدار عددی با تعداد یا نرخ اسپایک نمایش داده می‌شود:
  \[
  x=0.8,\quad T=4
  \quad\Longrightarrow\quad
  xT\approx3
  \quad\Longrightarrow\quad
  [1,1,1,0]
  \]
  ترتیب دقیق می‌تواند متفاوت باشد؛ تعداد اسپایک‌ها معنای اصلی را حمل می‌کند.
\end{itemize}

\subsection{کدگذاری مبتنی بر \eng{Integration}}

\begin{itemize}
  \item مقدار پیوسته در هر گام به پتانسیل اضافه می‌شود؛ با عبور از آستانه، اسپایک تولید و آستانه از پتانسیل کم می‌شود.
  \item برای \(x=0.8\) و \(\theta=1\):
  \begin{enumerate}
    \item \(0.8<1\): بدون اسپایک.
    \item \(0.8+0.8=1.6\): اسپایک؛ باقی‌مانده \(0.6\).
    \item \(0.6+0.8=1.4\): اسپایک؛ باقی‌مانده \(0.4\).
    \item \(0.4+0.8=1.2\): اسپایک.
  \end{enumerate}
  \[
  \boxed{[0,1,1,1]}
  \]
  \item حداقل شرط کدگذاری معنادار: \key{رابطهٔ مستقیم مقدار پیوسته با نرخ اسپایک}.
\end{itemize}

\begin{trap}{مبادلهٔ اصلی}
\[
\mathrm{Energy/Cost}\quad\leftrightarrow\quad\mathrm{Accuracy}
\]
اسپایکی‌کردن داده می‌تواند بخشی از اطلاعات پیوسته را از بین ببرد؛ بنابراین \eng{SNN} معمولاً انرژی کمتری مصرف می‌کند، اما ممکن است دقت پایین‌تری از \eng{Transformer} معمولی داشته باشد.
\end{trap}

\subsection{هدف \eng{Attention} اسپایکی}

\begin{itemize}
  \item مانند توجه معمولی، بخش‌های مرتبط با توجه به زمینه به یکدیگر وزن بیشتری می‌دهند.
  \item در متن، زمینهٔ «جنگل» باید بازنمایی «شیر» را به معنای حیوانی نزدیک کند؛ در تصویر نیز ناحیه‌های مرتبط برجسته می‌شوند.
  \item در فضای اسپایکی، ارتباط بیشتر معمولاً با \key{نرخ یا هم‌وقوعی بیشتر اسپایک‌ها} نشان داده می‌شود.
\end{itemize}

\subsection{\eng{Query}، \eng{Key} و \eng{Value} اسپایکی}

\begin{itemize}
  \item سه ماتریس وزن جداگانه داریم:
  \[
  \boxed{W_Q,\qquad W_K,\qquad W_V}
  \]
  \item دنبالهٔ اسپایکی ورودی در این وزن‌ها ضرب می‌شود و \(Q\)، \(K\) و \(V\) اسپایکی ساخته می‌شوند.
  \item وزن‌ها در آموزش با \eng{Surrogate Gradient} به‌روزرسانی می‌شوند.
\end{itemize}

\subsection{\eng{Attention} معمولی در برابر \eng{Spiking Attention}}

\begin{itemize}
  \item توجه کلاسیک:
  \[
  \boxed{\operatorname{softmax}\left(\frac{QK^T}{\sqrt d}\right)V}
  \]
  \eng{Softmax} نمره‌ها را به وزن‌هایی نرمال می‌کند که مجموعشان برابر یک است.
  \item توجه اسپایکی:
  \[
  \boxed{QK^TV}
  \]
  ضرب اسپایک‌های دودویی خود نوعی هم‌فعالی یا شباهت می‌سازد؛ بنابراین \eng{Softmax} ضرورت و معنای قبلی را ندارد. مقیاس‌گذاری با \(\sqrt d\) همچنان می‌تواند باقی بماند.
  \item حذف \eng{Softmax} اجازه می‌دهد ترتیب ضرب‌ها تغییر کند:
  \[
  \boxed{(QK^T)V=Q(K^TV)}
  \]
  ابتدا محاسبهٔ \(K^TV\) و سپس ضرب در \(Q\) می‌تواند پیچیدگی، حافظه و انرژی را کاهش دهد.
\end{itemize}

\subsection{معماری کلی \eng{SpikeFormer}}

\begin{enumerate}
  \item کدگذاری ورودی پیوسته به دنبالهٔ اسپایکی.
  \item ورود به \term{توجه خودی اسپایکی}{Spiking Self-Attention; SSA}.
  \item عبور خروجی توجه از بلوک \eng{MLP} اسپایکی.
  \item استفاده از اتصال‌های \eng{Residual} اطراف بلوک توجه و \eng{MLP}.
  \item تصمیم نهایی با سر طبقه‌بندی یا خروجی وظیفه.
\end{enumerate}

\subsection{نقش \eng{LIF} در \eng{SpikeFormer}}

\begin{itemize}
  \item پس از بسیاری از لایه‌های خطی، نورون \eng{LIF} قرار می‌گیرد.
  \item وظایف \eng{LIF}: تجمع زمانی پتانسیل، نشت، مقایسه با آستانه، تولید اسپایک، ریست و دورهٔ تحریک‌ناپذیری.
  \item \eng{MLP} معمولی:
  \[
  \mathrm{Linear}\rightarrow\mathrm{Activation}\rightarrow\mathrm{Linear}
  \]
  \item \eng{MLP} اسپایکی:
  \[
  \boxed{\mathrm{Linear}\rightarrow\mathrm{LIF}
  \rightarrow\mathrm{Linear}\rightarrow\mathrm{LIF}}
  \]
  \item \eng{LIF} معادل دقیق تابع فعال‌ساز نیست، اما غیرخطی‌بودن ایجاد و خروجی را دوباره اسپایکی می‌کند.
\end{itemize}

\subsection{\eng{Residual Connection}}

\begin{itemize}
  \item اطلاعات قبلی را از مسیر میان‌بُر حفظ می‌کند.
  \item به کاهش مشکل \eng{vanishing gradient} کمک می‌کند.
  \item اگر یک بلوک بازنمایی مفیدی نسازد، بازنمایی قبلی کاملاً از بین نمی‌رود.
\end{itemize}

\section{بخش دوم: \eng{SpikeBERT} و تقطیر دانش}

\subsection{چرا متن برای \eng{SNN} دشوارتر است؟}

\begin{itemize}
  \item \term{تراکم اطلاعات}{Information Density}: تغییر کوچک مانند افزودن «نـ» می‌تواند معنای جمله را معکوس کند؛ تصویر معمولاً افزونگی بیشتری دارد.
  \item \term{وابستگی بلندمدت}{Long-Term Dependency}: معنای کلمه ممکن است به بخش‌های بسیار دور متن وابسته باشد، اما نشت \eng{LIF} نوعی فراموشی ایجاد می‌کند.
  \item \term{ماهیت نمادین زبان}{Symbolic Nature}: توکن گسسته ابتدا به \eng{embedding} پیوسته و سپس دوباره به اسپایک تبدیل می‌شود؛ احتمال ازدست‌رفتن ظرافت معنایی افزایش می‌یابد.
  \item \eng{SpikeFormer} و \eng{SpikeBERT} بیشتر \term{اثبات مفهوم}{Proof of Concept} هستند؛ امکان تبدیل معماری را نشان می‌دهند، اما هنوز چالش دقت و همگرایی دارند.
\end{itemize}

\subsection{\eng{Knowledge Distillation}}

\begin{itemize}
  \item آموزش \eng{SpikeBERT} فقط با برچسب نهایی به‌خوبی همگرا نمی‌شود؛ بنابراین خروجی‌های میانی نیز راهنمایی می‌شوند.
  \item \textbf{\eng{Teacher}}: مدل \eng{BERT} معمولی، ازپیش‌آموزش‌دیده و دارای بازنمایی‌های پیوستهٔ زبانی.
  \item \textbf{\eng{Student}}: مدل اسپایکی \eng{SpikeBERT}.
  \item دانش از چهار مسیر منتقل می‌شود:
  \begin{itemize}
    \item \eng{Embedding Loss}
    \item \eng{Hidden Feature Loss}
    \item \eng{Logit Loss}
    \item \eng{Cross Entropy Loss} نسبت به برچسب واقعی
  \end{itemize}
\end{itemize}

\subsection{تبدیل \eng{Embedding} به اسپایک}

\begin{itemize}
  \item شکل ورودی پیوسته:
  \[
  \boxed{X\in\mathbb{R}^{N\times D}}
  \]
  \(N\): تعداد توکن‌ها؛ \(D\): تعداد ویژگی‌های هر توکن.
  \item با افزودن \(T\) گام زمانی:
  \[
  \boxed{N\times D\quad\longrightarrow\quad T\times N\times D}
  \]
  \item مثال با \(T=5\)، مقدار \(0.3\) و آستانهٔ \(1\):
  \[
  0.3\rightarrow0.6\rightarrow0.9\rightarrow1.2
  \rightarrow\text{\rl{اسپایک و ریست}}
  \]
\end{itemize}

\subsection{بازگردانی و هم‌ترازسازی بازنمایی}

\begin{enumerate}
  \item دریافت \eng{embedding} اولیه از \eng{BERT}.
  \item تبدیل مقدار پیوسته به اسپایک.
  \item پردازش در \eng{SpikeBERT}.
  \item تبدیل تعداد یا نرخ اسپایک به مقدار پیوسته.
  \item نگاشت خروجی \eng{Student} به فضای \eng{Teacher}.
  \item مقایسه و محاسبهٔ خطا.
\end{enumerate}

\begin{formula}{Projection}
\[
\boxed{H'=HW+b}
\qquad\text{\rl{یا}}\qquad
\boxed{H'=\mathrm{MLP}(H)}
\]
\eng{MLP} نسبت به یک لایهٔ خطی ساده، نگاشت قوی‌تری برای هم‌ترازکردن فضای دو مدل می‌سازد.
\end{formula}

\subsection{\eng{Feature Loss} و \eng{Embedding Loss}}

\begin{itemize}
  \item خروجی‌های متناظر \eng{Teacher} و \eng{Student} پس از \eng{projection} خانه‌به‌خانه مقایسه می‌شوند.
  \item مثال:
  \[
  H_T=[1,2,0],\qquad H_S=[0.5,0.3,1]
  \]
  \[
  \boxed{L=(1-0.5)^2+(2-0.3)^2+(0-1)^2}
  \]
  \item لازم نیست همهٔ لایه‌ها هم‌تراز شوند؛ انتخاب لایه‌های مقایسه‌شونده بخشی از طراحی مدل است.
\end{itemize}

\subsection{دو مرحلهٔ آموزش \eng{SpikeBERT}}

\begin{itemize}
  \item \textbf{مرحلهٔ اول ـ انتقال دانش عمومی:}
  \begin{itemize}
    \item نزدیک‌کردن \eng{embedding}ها و ویژگی‌های میانی \eng{Student} به \eng{Teacher}.
    \item تمرکز اصلی روی \eng{Embedding Loss} و \eng{Hidden Feature Loss}.
  \end{itemize}
  \item \textbf{مرحلهٔ دوم ـ آموزش مخصوص وظیفه:}
  \begin{itemize}
    \item تنظیم مدل برای وظیفه‌ای مانند تحلیل احساسات.
    \item مقایسهٔ \eng{logit}های دو مدل و خروجی \eng{Student} با برچسب واقعی.
  \end{itemize}
\end{itemize}

\subsection{\eng{Logit Loss} و \eng{Cross Entropy}}

\begin{itemize}
  \item مثال خطای لاجیت:
  \[
  z_T=[0.2,0.1],\qquad z_S=[0.5,0]
  \]
  \[
  \boxed{L_{\mathrm{logit}}=(0.2-0.5)^2+(0.1-0)^2}
  \]
  \item خطای آنتروپی متقاطع برای کلاس درست \(y\):
  \[
  \boxed{
  L_{\mathrm{CE}}=
  -\log\left(\frac{e^{z_y}}{\sum_j e^{z_j}}\right)}
  \]
  \(z_y\): لاجیت کلاس صحیح؛ مخرج: مجموع نمایی لاجیت همهٔ کلاس‌ها.
\end{itemize}

\subsection{تابع زیان نهایی و آموزش}

\begin{formula}{زیان ترکیبی SpikeBERT}
\[
\boxed{
\begin{aligned}
L_{\mathrm{total}}={}&
\lambda_1L_{\mathrm{embedding}}
+\lambda_2L_{\mathrm{feature}}\\
&+\lambda_3L_{\mathrm{logit}}
+\lambda_4L_{\mathrm{CE}}
\end{aligned}}
\]
\begin{itemize}
  \item \(\lambda_i\): وزن اهمیت هر مؤلفهٔ خطا.
  \item پس‌انتشار بر اساس \(L_{\mathrm{total}}\) انجام و برای عبور گرادیان از اسپایک‌ها از \eng{Surrogate Gradient} استفاده می‌شود.
  \item حتی \eng{embedding} اولیهٔ \eng{SpikeBERT} نیز در آموزش \eng{fine-tune} می‌شود.
\end{itemize}
\end{formula}

\subsection{نقش \eng{Time Step}}

\begin{itemize}
  \item \(T\) یک ابرپارامتر ویژهٔ شبکهٔ اسپایکی و تعداد گام‌های زمانی کدگذاری است:
  \[
  \boxed{T=\mathrm{Time\ Step}}
  \]
  \item افزایش \(T\) مانند \term{تجمع شواهد}{Evidence Accumulation} است؛ شواهد بیشتری برای عبور از آستانه و تصمیم‌گیری فراهم می‌شود.
  \item رابطهٔ دقت تا نقطه‌ای صعودی و سپس اشباعی است:
  \[
  \boxed{T\uparrow\quad\Longrightarrow\quad
  \mathrm{Accuracy}\uparrow\quad\text{\rl{تا نقطهٔ اشباع}}}
  \]
  \item هدف: یافتن کمترین \(T\) که بهترین دقت عملی یا \eng{temporal resolution} مناسب را ایجاد کند.
  \item مشکلات \(T\) بسیار بزرگ:
  \begin{itemize}
    \item پایان اطلاعات مفید و حساسیت بیشتر به نویز؛
    \item افزایش بار حافظه و ارتباط با محدودیت حافظهٔ فعال؛
    \item کندترشدن همگرایی، آموزش و \eng{inference}؛
    \item هزینهٔ بیشتر در برابر بهبود اندک، ثابت‌ماندن یا حتی افت دقت.
  \end{itemize}
\end{itemize}

\subsection{مثال \eng{Attention} اسپایکی: شیر و جنگل}

\begin{itemize}
  \item با فرض \(Q=K=V=X\)، امتیاز «شیر به شیر» برابر \(4\) و «شیر به جنگل» برابر \(2\) است.
  \[
  4+2=6,\qquad
  A_{\text{\rl{شیر، شیر}}}=\frac{4}{6}=\frac23,\qquad
  A_{\text{\rl{شیر، جنگل}}}=\frac{2}{6}=\frac13
  \]
  \item بازنمایی جدید:
  \[
  \boxed{
  V'_{\text{\rl{شیر}}}
  =\frac23V_{\text{\rl{شیر}}}
  +\frac13V_{\text{\rl{جنگل}}}}
  \]
  زمینهٔ «جنگل» باید بُعد حیوانی «شیر» را تقویت و بُعد مایع را تضعیف کند.
\end{itemize}

\begin{trap}{جمع‌بندی مقایسه‌ای}
\begin{itemize}
  \item \textbf{\eng{SpikeFormer}:} تمرکز بر معماری \eng{Transformer} اسپایکی، کدگذاری، توجه بدون \eng{Softmax} و \eng{MLP} مبتنی بر \eng{LIF}.
  \item \textbf{\eng{SpikeBERT}:} انتقال همین ایده به زبان با کمک \eng{Teacher--Student} و چند زیان میانی و نهایی.
  \item \textbf{مبادلهٔ مشترک:} کاهش انرژی و پیچیدگی در برابر خطر افت دقت و ازدست‌رفتن اطلاعات پیوسته.
\end{itemize}
\end{trap}

\section{بخش سوم: \eng{SpikeGPT}، \eng{RWKV} و تولید متن اسپایکی}

\subsection{طبقه‌بندی در برابر تولید}

\begin{itemize}
  \item \term{طبقه‌بندی}{Classification}:
  \begin{itemize}
    \item \eng{BERT} دوطرفه و \eng{encoder-only} است؛ برای تصمیم نهایی کل متن را از دو جهت می‌بیند.
    \item \eng{hidden state}ها می‌توانند \eng{average pooling} شوند و سپس به سر طبقه‌بندی برسند.
  \end{itemize}
  \item \term{تولید}{Generation}:
  \begin{itemize}
    \item \eng{GPT} علّی، خودرگرسیو و \eng{decoder-only} است.
    \item با \eng{masked self-attention} فقط گذشته را می‌بیند و توکن بعدی را پیش‌بینی می‌کند.
    \item در آموزش، هدف‌ها یک موقعیت شیفت می‌یابند:
    \[
    [x_1,x_2,\ldots,x_{s-1}]
    \longrightarrow
    [x_2,x_3,\ldots,x_s]
    \]
    \item خروجی به واژگان متصل و معمولاً با \eng{Cross Entropy} آموزش داده می‌شود.
  \end{itemize}
\end{itemize}

\subsection{چرا \eng{SpikeGPT} از \eng{RWKV} استفاده می‌کند؟}

\begin{itemize}
  \item تبدیل مستقیم \eng{Transformer} به مدل اسپایکی برای تولید متن معمولاً کافی و کارآمد نیست.
  \item توجه کامل همهٔ توکن‌ها را دوبه‌دو مقایسه می‌کند و پیچیدگی تقریبی دارد:
  \[
  \boxed{O(n^2)}
  \]
  \item \eng{RNN} گذشته را در یک \eng{state} خلاصه می‌کند، اما غیرخطی‌بودن‌هایی مانند \(\tanh\) روی حالت، فراموشی و \eng{vanishing gradient} ایجاد می‌کنند.
  \item \textbf{\eng{RWKV}}: ترکیب حالت بازگشتی \eng{RNN} با مزایای \eng{Transformer}؛ حالت خطی‌تر است و گذشته بدون توجه دوبه‌دوی کامل خلاصه می‌شود.
\end{itemize}

\subsection{چهار جزء \eng{RWKV}}

\begin{itemize}
  \item \term{\eng{Receptance}}{R}: گیت خروجی؛ تعیین می‌کند چه مقدار از \eng{state} عبور کند.
  \item \term{\eng{Weight}}{W}: عامل فراموشی؛ مقدار حفظ یا کاهش اطلاعات قبلی.
  \item \term{\eng{Key}}{K}: اهمیت توکن فعلی.
  \item \term{\eng{Value}}{V}: محتوای توکن فعلی.
  \item در \eng{RWKV}، \eng{Query} حذف می‌شود؛ گذشته در یک \eng{state} جمع شده است.
\end{itemize}

\begin{formula}{خروجی \eng{RWKV}}
\[
\boxed{\mathrm{Output}=R\times\mathrm{State}}
\]
\end{formula}

\subsection{مثال عددی به‌روزرسانی \eng{State}}

\begin{enumerate}
  \item توکن اول:
  \[
  K_1=2,\quad V_1=10,\quad
  \mathrm{State}_1=\frac{K_1V_1}{K_1}=\frac{20}{2}=10
  \]
  \item توکن دوم و عامل فراموشی:
  \[
  K_2=3,\quad V_2=4,\quad R_2=0.9,\quad W=0.5
  \]
  \item صورت و مخرج جدید:
  \[
  0.5(20)+3(4)=22,\qquad 0.5(2)+3=4
  \]
  \item حالت و خروجی:
  \[
  \boxed{\mathrm{State}_2=\frac{22}{4}=5.5}
  \qquad
  \boxed{\mathrm{Output}=0.9(5.5)=4.95}
  \]
\end{enumerate}

\begin{trap}{دو معنای متفاوت زمان}
\begin{itemize}
  \item \(t\) در \eng{RWKV}: موقعیت توکن در دنبالهٔ زبانی.
  \item \(T\) در \eng{SNN}: تعداد گام‌های زمانی برای کدگذاری اسپایکی هر نمایش.
\end{itemize}
\eng{SpikeGPT} هم‌زمان با زمان ترتیبی زبان و زمان اسپایکی سروکار دارد.
\end{trap}

\subsection{روند معماری \eng{SpikeGPT}}

\begin{enumerate}
  \item تبدیل \eng{embedding} ورودی به نمایش باینری یا اسپایکی با \eng{LIF}.
  \item ورود دادهٔ اسپایکی به بلوک \eng{RWKV}.
  \item اسپایکی‌کردن دوبارهٔ خروجی \eng{RWKV}.
  \item ورود به بخش فیدفوروارد اصلاح‌شده یا \eng{R-FFN}.
  \item اسپایکی‌کردن دوبارهٔ خروجی \eng{R-FFN}.
  \item استفاده از خروجی برای تولید متن یا طبقه‌بندی.
\end{enumerate}

\subsection{\eng{Time-Mix} و \eng{Channel-Mix}}

\begin{itemize}
  \item \term{\eng{Time-Mix}}{RWKV}: اطلاعات توکن‌های قبلی و توکن فعلی را ترکیب می‌کند؛ معادل نقش فهم زمینه در \eng{Attention}.
  \item \term{\eng{Channel-Mix}}{R-FFN}: روی ویژگی‌ها و کانال‌های توکن فعلی کار می‌کند؛ معادل نقش \eng{MLP} یا \eng{feed-forward}.
  \item در \eng{R-FFN}، بُعد نمایش ابتدا بزرگ، سپس از \eng{ReLU} عبور و دوباره فشرده می‌شود؛ مسیر گیت با محتوای اصلی عنصر‌به‌عنصر ترکیب می‌شود.
\end{itemize}

\begin{formula}{ضرب هادامارد در \eng{R-FFN}}
\[
\boxed{
[a_1,a_2,a_3]\odot[b_1,b_2,b_3]
=[a_1b_1,a_2b_2,a_3b_3]}
\]
\end{formula}

\subsection{پیوستگی میانی، تنکی و مصرف انرژی}

\begin{itemize}
  \item لازم نیست همهٔ نقاط \eng{SNN} صفر و یک باشند؛ مقدارهای پیوستهٔ میانی عمدتاً پتانسیل غشا هستند.
  \item خروجی لایه‌های اسپایکی بسیار \term{تنک}{Sparse} است؛ فقط رخدادهای فعال محاسبه ایجاد می‌کنند.
  \item \eng{GPU} برای عملیات متراکم طراحی شده است؛ مزیت واقعی \eng{SNN} روی سخت‌افزار نورومورفیک آشکارتر می‌شود.
  \item عملیات شبکهٔ غیر اسپایکی:
  \[
  \boxed{\mathrm{MAC}=\mathrm{Multiply}+\mathrm{Accumulate}}
  \]
  \item با ورودی صفر و یک در شبکهٔ اسپایکی، بسیاری از ضرب‌ها حذف می‌شوند:
  \[
  \boxed{\mathrm{AC}=\mathrm{Accumulate}}
  \]
  \item برآورد مقالهٔ \eng{SpikeGPT}: کاهش انرژی تقریبی \key{۳۲ برابری} نسبت به \eng{Vanilla GPT}.
\end{itemize}

\begin{trap}{جمع‌بندی فصل}
\begin{itemize}
  \item \textbf{\eng{SpikeFormer}}: توجه اسپایکی کم‌هزینه‌تر برای معماری‌های عمیق.
  \item \textbf{\eng{SpikeBERT}}: حل دشواری متن با تقطیر دانش از \eng{BERT}.
  \item \textbf{\eng{SpikeGPT}}: تولید خودرگرسیو با \eng{RWKV}، \eng{Time-Mix} و \eng{R-FFN} به‌جای توجه کامل.
\end{itemize}
\end{trap}

\end{document}
